\chapter{Solutions: Final Exam, June 25, 2015}
\label{s:physics-2015-06-final}
The electronic and vibrational conventions are stated explicitly below. In particular, a band is counted before separately resolving spin unless otherwise specified, and phonon polarization vectors are mass-weighted as in the hint.

\section*{1. Bloch sums and their role in tight binding}
Let $\phi_{\nu a}(\vct r-\vct R-\vct\tau_\nu)$ be a localized orbital of type $a$ on basis atom $\nu$ in the cell at $\vct R$. For $N$ cells with periodic boundary conditions, define
\[
\boxed{\Phi_{\nu a,\vct k}(\vct r)
=\frac1{\sqrt N}\sum_{\vct R}\ee^{\ii\vct k\cdot\vct R}
\phi_{\nu a}(\vct r-\vct R-\vct\tau_\nu)}.
\]
The factor $N^{-1/2}$ is conventional; if orbitals on different sites overlap, the overlap matrix still determines the final normalization.

For any Bravais translation $\vct T$, reindex the sum by $\vct R'=\vct R-\vct T$:
\begin{align*}
\Phi_{\nu a,\vct k}(\vct r+\vct T)
&=\frac1{\sqrt N}\sum_{\vct R'}
\ee^{\ii\vct k\cdot(\vct R'+\vct T)}
\phi_{\nu a}(\vct r-\vct R'-\vct\tau_\nu)\\
&=\ee^{\ii\vct k\cdot\vct T}\Phi_{\nu a,\vct k}(\vct r).
\end{align*}
The sum therefore has precisely the Bloch translation property.

In tight binding, these symmetry-adapted localized-orbital sums form the basis at each wave vector:
\[
\psi_{n\vct k}=\sum_{\nu a}c_{\nu a}^{(n)}(\vct k)\Phi_{\nu a,\vct k}.
\]
Projecting the Schr\"odinger equation onto this basis gives
\[
\boxed{H(\vct k)c_n=E_n(\vct k)S(\vct k)c_n},
\]
where $H_{ij}=\langle\Phi_i|\hat H|\Phi_j\rangle$ and
$S_{ij}=\langle\Phi_i|\Phi_j\rangle$. In an orthonormal model, $S=I$.
Translational symmetry has reduced an extensive real-space problem to a finite matrix at each $\vct k$; hopping between localized orbitals determines the matrix entries and their wave-vector-dependent phases.

\section*{2. Five orbitals per Si atom and the band count}
For electrical and bonding properties, retain the valence $3s$ and three $3p$ orbitals. A useful fifth orbital is an excited effective $s$-like orbital $s^*$, giving the conventional
\[
\boxed{\{3s,3p_x,3p_y,3p_z,s^*\}\quad(sp^3s^*)}
\]
basis. The $s^*$ orbital improves the description of conduction bands; in an empirical model it is an effective excited orbital, not a claim that the isolated atom already occupies a fifth valence orbital. The filled $1s$, $2s$, and $2p$ shells are treated as core states for this low-energy model. The original model is described by
\href{https://www.sciencedirect.com/science/article/pii/0022369783900641}{Vogl, Hjalmarson, and Dow (1983)}.

Silicon has two atoms per diamond primitive cell. Five orbital functions per atom therefore yield
\[
\boxed{2\times5=10\ \text{bands before explicit spin resolution}}.
\]
Each atom supplies four valence electrons, so each primitive cell supplies eight. A filled band accommodates two electrons per primitive cell when both spins are included. Hence
\[
\boxed{\frac{2\times4}{2}=4\ \text{fully occupied valence bands}}.
\]
The other six bands of this model are unoccupied in the undoped ground state. If spin is represented explicitly, the matrix has dimension twenty and eight spin-resolved bands are occupied. These conventions are consistent with the
\href{https://www.nextnano.com/docu/nextnano3/tutorials/1D_TightBinding_bulk.html}{nextnano implementation of the $sp^3s^*$ model}.

\section*{3. Why a direct all-electron plane-wave expansion is expensive}
At fixed $\vct k$, a Bloch wave can be expanded as
\[
\psi_{n\vct k}(\vct r)=\sum_{\vct G}c_{n\vct k}(\vct G)
\ee^{\ii(\vct k+\vct G)\cdot\vct r}.
\]
Near a nucleus, the all-electron potential varies strongly and the wavefunctions contain short spatial scales. Valence functions must also remain orthogonal to localized core states, producing nodes and rapid core-region variation. Resolving a length scale $\ell$ requires components with $|\vct k+\vct G|$ of order $1/\ell$. A kinetic-energy cutoff
\[
\frac{\hbar^2|\vct k+\vct G|^2}{2m_e}\le E_{\rm cut}
\]
then requires a very large $E_{\rm cut}$ and many plane waves, even if only valence energies are sought. The basis is valid, but direct all-electron convergence can be inefficient.

One remedy is to replace the core-region problem by a pseudopotential with smooth pseudo-valence functions. Alternatively, augmented-plane-wave or projector-augmented-wave approaches retain or reconstruct atomic core-region structure while using smoother representations elsewhere. Orthogonalized-plane-wave methods also account explicitly for the core orthogonality constraint. In each case, the objective is to represent the sharp atomic behavior without forcing a uniform plane-wave basis to resolve every core length scale directly.

\section*{4. Pseudopotentials and smooth pseudo-wavefunctions}
A pseudopotential is an effective electron--ion operator replacing the explicit nuclear attraction and core-electron effects in a valence calculation. It is constructed to reproduce relevant valence eigenvalues and scattering behavior outside a chosen core region. It can be angular-momentum dependent and nonlocal; it need not be a single scalar function of position.

A pseudo-wavefunction follows the required valence behavior outside the core but replaces the rapid all-electron oscillations inside it by a smoother function. Its large-$|\vct G|$ Fourier coefficients are therefore much smaller. A substantially lower plane-wave cutoff can represent it to a chosen accuracy.

This is an approximation with a chosen core--valence partition and transferability range. Smoothness does not itself guarantee correct bonding in every environment, and the cutoff must still be converged. Practical construction and validation requirements are discussed in
\href{https://www.quantum-espresso.org/wp-content/uploads/2022/03/pseudo-gen.pdf}{Quantum ESPRESSO, Notes on pseudopotential generation}.

\section*{5. The adiabatic energy surface}
Collect all nuclear positions into $\vct R$ and the electronic coordinates into $\vct r$. For fixed nuclei, solve the electronic eigenproblem
\[
\hat H_e(\vct R)\Psi_n(\vct r;\vct R)
=\mathcal E_n(\vct R)\Psi_n(\vct r;\vct R).
\]
The adiabatic potential-energy surface for nuclear motion is
\[
\boxed{E_n(\vct R)=\mathcal E_n(\vct R)+V_{NN}(\vct R)},
\]
where $V_{NN}$ is the nucleus--nucleus repulsion. If that term has already been included in $\hat H_e$, it must not be added a second time.

The electronic eigenproblem is solved parametrically as the nuclei move. In the Born--Oppenheimer approximation, nuclei move on a selected surface, neglecting transitions to other electronic surfaces and the associated nonadiabatic couplings. Nuclear kinetic energy is not part of this potential-energy surface. Its gradient determines the forces, and its Hessian determines harmonic restoring forces about equilibrium.

\section*{6. Obtaining lattice normal modes from the ground surface}
\subsection*{Equilibrium and harmonic force constants}
First locate a stationary nuclear configuration $\vct R^0$ with
$\partial E_0/\partial R_{i\alpha}=0$. Let
$u_{i\alpha}=R_{i\alpha}-R^0_{i\alpha}$, where $i$ labels atoms and $\alpha$ labels Cartesian components. Expand
\[
E_0(\vct R^0+\vct u)=E_0(\vct R^0)
+\frac12\sum_{i\alpha,j\beta}\Phi_{i\alpha,j\beta}
 u_{i\alpha}u_{j\beta}+O(u^3),
\qquad
\Phi_{i\alpha,j\beta}
=\left.\frac{\partial^2E_0}{\partial R_{i\alpha}\partial R_{j\beta}}\right|_{\vct R^0}.
\]
The linear term vanishes at equilibrium. The Hessian is the harmonic force-constant matrix, obtained analytically or by converged finite differences of energies or forces.

\subsection*{Mass weighting and the eigenproblem}
Newton's equations in the harmonic approximation are
\[
M_i\ddot u_{i\alpha}=-\sum_{j\beta}\Phi_{i\alpha,j\beta}u_{j\beta}.
\]
Substitute $u_{i\alpha}(t)=A_{i\alpha}\ee^{-\ii\omega t}/\sqrt{M_i}$. This gives
\[
\sum_{j\beta}\mathcal D_{i\alpha,j\beta}A_{j\beta}=\omega^2 A_{i\alpha},
\qquad
\boxed{\mathcal D_{i\alpha,j\beta}
=\frac{\Phi_{i\alpha,j\beta}}{\sqrt{M_iM_j}}}.
\]
Diagonalizing the real-space mass-weighted Hessian yields the squared normal-mode frequencies and orthonormal polarization vectors. The physical displacements are obtained by dividing each atomic block of an eigenvector by $\sqrt{M_i}$.

\subsection*{Periodic crystal and stability}
For a periodic crystal with $r$ atoms per primitive cell, Fourier transformation block-diagonalizes the problem into a $3r\times3r$ dynamical matrix at each allowed $\vct q$. The next answer gives its entries. The branches $\omega_p(\vct q)$ and their eigenvectors are the lattice normal modes.

Translational invariance imposes the acoustic sum rule
\[
\sum_j\Phi_{i\alpha,j\beta}=0
\]
for every $i,\alpha,\beta$. Thus three rigid translations have zero frequency at $\vct q=0$. A mechanically stable harmonic equilibrium has no negative $\omega^2$; negative eigenvalues describe an instability rather than ordinary oscillations. A nondegenerate, well-separated \emph{electronic} ground surface justifies the adiabatic treatment but does not by itself prove that a chosen nuclear configuration is a local minimum.

\section*{7. Dynamical-matrix elements and phase convention}
To remove a notational ambiguity, use the symbol $D$ exactly as in the question for the force constants, and $\mathcal D$ for the mass-weighted dynamical matrix. Adopt
\[
D_{\nu\alpha,\nu'\alpha'}(\vct t_n)
=\left.\frac{\partial^2E_0}
{\partial u_{0\nu\alpha}\,\partial u_{n\nu'\alpha'}}\right|_0.
\]
Choose the displacement convention
\[
u_{n\nu\alpha}(t)=\frac{A_{\nu\alpha}(\vct q)}{\sqrt{M_\nu}}
\ee^{\ii\vct q\cdot\vct t_n-\ii\omega t}.
\]
Inserting this into the equation of motion for atom $(0,\nu)$ gives
\[
\boxed{\mathcal D_{\nu\alpha,\nu'\alpha'}(\vct q)
=\frac1{\sqrt{M_\nu M_{\nu'}}}
\sum_n D_{\nu\alpha,\nu'\alpha'}(\vct t_n)
\ee^{\ii\vct q\cdot\vct t_n}},
\]
and the eigenproblem
\[
\sum_{\nu'\alpha'}\mathcal D_{\nu\alpha,\nu'\alpha'}(\vct q)
A_{\nu'\alpha'}=\omega^2 A_{\nu\alpha}.
\]
Hessian symmetry and translation invariance give
$D_{\nu\alpha,\nu'\alpha'}(\vct t_n)
=D_{\nu'\alpha',\nu\alpha}(-\vct t_n)$, so $\mathcal D(\vct q)$ is Hermitian.

If the plane-wave factor is instead defined using the full equilibrium position
$\vct t_n+\vct\tau_\nu$, the exponential in the boxed matrix becomes
$\ee^{\ii\vct q\cdot(\vct t_n+\vct\tau_{\nu'}-\vct\tau_\nu)}$.
The two matrices are related by a diagonal unitary change of basis and have identical frequencies. Stating the convention prevents an apparent phase disagreement from becoming a physical one.

\section*{8. Counting BaTiO$_3$ vibrational degrees of freedom}
There are $r=5$ atoms per cell and three displacement coordinates per atom.
\begin{enumerate}[label=(\alph*)]
\item At each $\vct q$, the Fourier-space dynamical matrix acts on $3r=15$ coordinates, so its dimension is
\[
\boxed{15\times15}.
\]
If ``dynamical matrix'' instead refers to the complete real-space system of $N$ cells, its dimension is $15N\times15N$. The standard band/branch formulation uses the first convention.
\item There are $N$ allowed wave vectors and fifteen modes at each one:
\[
\boxed{15N\ \text{normal modes}}.
\]
The real-displacement condition relates complex amplitudes at $\vct q$ and $-\vct q$; it does not halve the total number of real mechanical degrees of freedom.
\item The three rigid translations evolve into three acoustic branches:
\[
\boxed{3\ \text{acoustic branches}}.
\]
\item The remaining branches are optical:
\[
\boxed{3r-3=12\ \text{optical branches}}.
\]
Each branch contains $N$ normal modes, so the corresponding mode counts are $3N$ acoustic and $12N$ optical.
\end{enumerate}

\section*{9. Long-wavelength optical polarization vectors}
The hint uses \emph{mass-weighted} vectors $\vct A_\nu$, not the actual atomic displacements. At $\vct q=0$, a rigid translation in direction $\hat{\vct e}_\alpha$ has a mass-weighted eigenvector proportional to
\[
\left(\sqrt{M_1}\hat{\vct e}_\alpha,
      \sqrt{M_2}\hat{\vct e}_\alpha\right).
\]
A finite-frequency optical eigenvector of the Hermitian dynamical matrix is orthogonal to each of the three zero-frequency translational vectors. Component by component this implies
\[
\sqrt{M_1}\vct A_1+\sqrt{M_2}\vct A_2=0,
\qquad
\vct A_2=-\sqrt{\frac{M_1}{M_2}}\vct A_1.
\]
Therefore
\[
\vct A_1^*\cdot\vct A_2
=-\sqrt{\frac{M_1}{M_2}}\abs{\vct A_1}^2,
\qquad
\boxed{\frac{\abs{\vct A_1}^2}{\vct A_1^*\cdot\vct A_2}
=-\sqrt{\frac{M_2}{M_1}}}.
\]
The result is unchanged by a common complex phase of the mode. In terms of the \emph{physical} displacement amplitudes $\vct u_\nu=\vct A_\nu/\sqrt{M_\nu}$, the corresponding relation is
\[
\vct u_2=-\frac{M_1}{M_2}\vct u_1,
\qquad M_1\vct u_1+M_2\vct u_2=0.
\]
The unit-cell center of mass is stationary in the optical limit. The square-root mass ratio in the requested answer arises specifically from the mass-weighted convention in the hint.

\begin{center}
\begin{tikzpicture}[>=Latex]
\draw[gray,dashed] (-.5,0)--(5,0);
\fill[blue!60!black] (1,0) circle (5pt);
\fill[blue!60!black] (3.5,0) circle (7pt);
\draw[->,thick] (1,.4)--(2.2,.4) node[midway,above] {$\vct u_1$};
\draw[->,thick] (3.5,.4)--(2.8,.4) node[midway,above] {$\vct u_2$};
\node[below] at (1,-.1) {$M_1$};
\node[below] at (3.5,-.1) {$M_2$};
\node[below] at (2.25,-.65) {$M_1\vct u_1+M_2\vct u_2=0$};
\end{tikzpicture}
\end{center}
The schematic shows opposite motion; arrow lengths are illustrative. At an exactly soft optical instability, the optical subspace must be selected orthogonally to translations before taking the zero-frequency limit. The standard finite-frequency optical limit is assumed in the question.

\section*{References and checks}
The preceding formulas follow from the displayed translation, state-counting, and Hessian arguments. The primary sources used to check conventions are:
\begin{itemize}
\item \href{https://www.sciencedirect.com/science/article/pii/0022369783900641}{P. Vogl, H. P. Hjalmarson, and J. D. Dow, A Semi-empirical tight-binding theory of the electronic structure of semiconductors (1983)}: the $sp^3s^*$ basis.
\item \href{https://www.nextnano.com/docu/nextnano3/tutorials/1D_TightBinding_bulk.html}{nextnano, Empirical tight-binding $sp^3s^*$ band structure}: ten-orbital and twenty-spin-orbital matrix conventions.
\item \href{https://www.quantum-espresso.org/wp-content/uploads/2022/03/pseudo-gen.pdf}{Quantum ESPRESSO, Notes on pseudopotential generation}: effective-core construction and validation.
\item \href{https://www.quantum-espresso.org/Doc/INPUT_DYNMAT.html}{Quantum ESPRESSO, dynmat.x documentation}: orthogonal dynamical-matrix eigenvectors versus physical displacements divided by square roots of masses.
\item \href{https://davidtong.org/pdfs/teaching/solid-state-physics/solidstate4.pdf}{David Tong, Solid State Physics, Chapter 4}: lattice vibrations and acoustic/optical branches.
\end{itemize}
